%%%%%%%%%%%%%%%%%%%%%%%%%%%%%%%%%%%%%%%%%
% University Assignment Title Page 
% LaTeX Template
% Version 1.0 (27/12/12)
%
% This template has been downloaded from:
% http://www.LaTeXTemplates.com
%
% Original author:
% WikiBooks (http://en.wikibooks.org/wiki/LaTeX/Title_Creation)
%
% License:
% CC BY-NC-SA 3.0 (http://creativecommons.org/licenses/by-nc-sa/3.0/)
% 
% Instructions for using this template:
% This title page is capable of being compiled as is. This is not useful for 
% including it in another document. To do this, you have two options: 
%
% 1) Copy/paste everything between \begin{document} and \end{document} 
% starting at \begin{titlepage} and paste this into another LaTeX file where you 
% want your title page.
% OR
% 2) Remove everything outside the \begin{titlepage} and \end{titlepage} and 
% move this file to the same directory as the LaTeX file you wish to add it to. 
% Then add \input{./title_page_1.tex} to your LaTeX file where you want your
% title page.
%
% template taken from https://www.overleaf.com/latex/examples/title-page-with-logo/hrskypjpkrpd#.WDXHQKIrLdc
%%%%%%%%%%%%%%%%%%%%%%%%%%%%%%%%%%%%%%%%%
%\title{Title page with logo}
%----------------------------------------------------------------------------------------
%   PACKAGES AND OTHER DOCUMENT CONFIGURATIONS
%----------------------------------------------------------------------------------------

\documentclass[12pt]{article}
\usepackage[english]{babel}
\usepackage[utf8x]{inputenc}
\usepackage{algpseudocode}
\usepackage{algorithm}
\usepackage{algorithmicx}
\usepackage{amsmath}
\usepackage{courier}
\usepackage{graphicx}
\usepackage{hyperref}
\usepackage{listings}
\usepackage{url}
\usepackage[colorinlistoftodos]{todonotes}

\begin{document}

\begin{titlepage}

\newcommand{\HRule}{\rule{\linewidth}{0.5mm}} % Defines a new command for the horizontal lines, change thickness here

\center % Center everything on the page
 
%----------------------------------------------------------------------------------------
%   HEADING SECTIONS
%----------------------------------------------------------------------------------------

\textsc{\LARGE The University of Texas at Austin}\\[1.5cm] % Name of your university/college
\textsc{\Large Software Testing}\\[0.5cm] % Major heading such as course name
\textsc{\large Software Engineering - Option III}\\[0.5cm] % Minor heading such as course title

%----------------------------------------------------------------------------------------
%   TITLE SECTION
%----------------------------------------------------------------------------------------

\HRule \\[0.4cm]
{ \huge \bfseries Testing Distributed Systems - Kafka}\\[0.4cm] % Title of your document
\HRule \\[1.5cm]
 
%----------------------------------------------------------------------------------------
%   AUTHOR SECTION
%----------------------------------------------------------------------------------------

\begin{minipage}{0.4\textwidth}
\begin{flushleft} \large
\emph{Authors:}\\
Howie \textsc{Benefiel}\\
Dale \textsc{Pedinski} % Your name
\end{flushleft}
\end{minipage}
~
\begin{minipage}{0.4\textwidth}
\begin{flushright} \large
\emph{Professor:} \\
Dr. Sarfraz \textsc{Khurshid} % Supervisor's Name
\end{flushright}
\end{minipage}\\[2cm]

% If you don't want a supervisor, uncomment the two lines below and remove the section above
%\Large \emph{Author:}\\
%John \textsc{Smith}\\[3cm] % Your name

%----------------------------------------------------------------------------------------
%   DATE SECTION
%----------------------------------------------------------------------------------------

{\large \today}\\[2cm] % Date, change the \today to a set date if you want to be precise

%----------------------------------------------------------------------------------------
%   LOGO SECTION
%----------------------------------------------------------------------------------------

\includegraphics{logo.png}\\[1cm] % Include a department/university logo - this will require the graphicx package
 
%----------------------------------------------------------------------------------------

\vfill % Fill the rest of the page with whitespace

\end{titlepage}


\begin{abstract}
\end{abstract}

\section{Introduction}

\section{Kafka Architecture}

\section{Prior Work}

From our literature review, we found that the field of testing distributed systems is relatively unexplored.
This is backed up by the findings of Yuan et al. \cite{simpletesting}.

The authors of Yuan et al. analyzed a number of open-source distributed systems running production workloads in industry currently.
Their findings showed that a large percentage of catastrophic failures could have been prevented by simple tests and static analysis.

In Yuan et al.\'s paper, they studied the catastrophuc failures 5 distributed, data-intensive systems, Cassandra, HBase, Hadoop Distrbibuted File System (HDFS), Hadoop MapReduce, and Redis.
They defined a catastrophic failure as a failure which would result in unavailability of the system or data loss.
To determine the faults of the system, the authors analyzed failures reported on the respective system's bug tracker.
For each system, the authors analyzed the catastrophic failures.

In general, Yuan et al. found that the failure modes of distributed systems were complex.
The authors determined that the complexity of these failure modes was caused by the relatively large state space of distributed systems.
Quantitatively, they found that 77\% of failures required more thatn one event to produce the failure, and 90\% required no more than three.

The authors of Yuan et al. then analyzed the types of inputs required to produce the failure.
The most unexpected finding from this analysis is that 24\% of failures were caused by a node becoming unreachable.
This is surprising because a core tenant of distributed systems is fault-tolerance, i.e., the system should be able to handle node\(s\) becoming unavailable.

The authors also analyzed the determinism of failures.
In distributed systems, there is the possibility that a sequence of events could produce a failure sometimes while not producing a failure other times.
The authors found that 74\% of failures are deterministic.
This is a promising result because it means the failures are not by determined by the timing of the input events.

By far the most unexpceted result of this paper

\section{Duck Tape}

\section{Custom Testing Apparatus}



\newpage


\bibliographystyle{plain}
\bibliography{references}



\appendix

\section{Github Code}

\href{https://github.com/howinator/software-testing-kafka-project}{https://github.com/howinator/software-testing-kafka-project}

%\section{Thread-Local Example}
%\label{sec:tls}
%\lstinputlisting[language=Java, firstline=366, lastline=382]{../../java-bag/src/SwedishBag.java}
% \section{Some \LaTeX{} Examples}
% \label{sec:examples}

% \subsection{Sections}

% Use section and subse{}ction commands to organize your document. \LaTeX{} handles all the formatting and numbering automatically. Use ref and label commands for cross-references.

% \subsection{Comments}

% Comments can be added to the margins of the document using the \todo{Here's a comment in the margin!} todo command, as shown in the example on the right. You can also add inline comments too:

% \todo[inline, color=green!40]{This is an inline comment.}

% \subsection{Tables and Figures}

% Use the table and tabular commands for basic tables --- see Table~\ref{tab:widgets}, for example. You can upload a figure (JPEG, PNG or PDF) using the files menu. To include it in your document, use the includegraphics command as in the code for Figure~\ref{fig:frog} below.

% % Commands to include a figure:
% \begin{figure}
% \centering
% \includegraphics[width=0.5\textwidth]{frog.jpg}
% \caption{\label{fig:frog}This is a figure caption.}
% \end{figure}

% \begin{table}
% \centering
% \begin{tabular}{l|r}
% Item & Quantity \\\hline
% Widgets & 42 \\
% Gadgets & 13
% \end{tabular}
% \caption{\label{tab:widgets}An example table.}
% \end{table}

% \subsection{Mathematics}

% \LaTeX{} is great at typesetting mathematics. Let $X_1, X_2, \ldots, X_n$ be a sequence of independent and identically distributed random variables with $\text{E}[X_i] = \mu$ and $\text{Var}[X_i] = \sigma^2 < \infty$, and let
% $$S_n = \frac{X_1 + X_2 + \cdots + X_n}{n}
%       = \frac{1}{n}\sum_{i}^{n} X_i$$
% denote their mean. Then as $n$ approaches infinity, the random variables $\sqrt{n}(S_n - \mu)$ converge in distribution to a normal $\mathcal{N}(0, \sigma^2)$.

% \subsection{Lists}

% You can make lists with automatic numbering \dots

% \begin{enumerate}
% \item Like this,
% \item and like this.
% \end{enumerate}
% \dots or bullet points \dots
% \begin{itemize}
% \item Like this,
% \item and like this.
% \end{itemize}

% We hope you find write\LaTeX\ useful, and please let us know if you have any feedback using the help menu above.



\end{document}
